\section{Discussion}
\label{sec:discussion}

\subsection{Assembling the answer}

We can now answer the question posed in the title.

Rotating fluids make polygons because the regular $N$-gon occupies a
distinguished position in the M\"obius parameter space $\Cstar$: it sits
at $|\lam| = 1$, where spiral orbits become closed polygonal orbits
(Step~1, the conjugacy theorem).  This position is an energy maximum
among ring configurations of $N$ point vortices (Step~3, $H''(0) < 0$
for all $N \ge 3$).  In two-dimensional turbulence at negative
temperature---the regime of large-scale vortex condensation
\citep{Onsager1949}---systems are driven toward energy maxima (Step~4).
The regular polygon is therefore the most probable macrostate.
Atmospheric turbulence provides the thermostat that enables this
statistical relaxation (Step~5).  Once the system arrives at
$|\lam| \approx 1$, eigenvalue branch disconnection (within the WKB
approximation) prevents smooth perturbations of the jet speed from
moving the system to the hyperbolic branch (Step~2).

\emph{That} a polygon forms is universal---a consequence of geometry
(the existence of the $|\lam| = 1$ locus), statistics (negative
temperature selects energy maxima), and eigenvalue structure (the polygon
branch is disconnected from the radial branch within WKB).

\emph{Which} polygon forms is planet-specific.  Saturn's Rossby
dispersion relation gives $n^* = 5.62 \to 6$; Jupiter's vortex dynamics
selects $N = 5$ at the south pole and $N = 8$ at the north.  Different
mechanisms arrive at $|\lam| \approx 1$ through different doors---Rossby
stationarity for Saturn, Thomson stability for Jupiter---but they arrive
at the same geometric destination.

\subsection{What this framework does and does not explain}

\paragraph{Explained.}
Why polygons persist (branch disconnection within WKB).  Why the polygon is
statistically favored (Onsager negative temperature).  Why Saturn's
hexagon has six sides (Rossby wavenumber selection).  Why the amplitude
is $\varepsilon = 0.112$ (matched asymptotic expansion with derived $C$).
Why the hexagon--NPS interaction has a 62-year cycle (near-resonance).
Why Jupiter's cyclone rings are nearly perfect polygons
($\sigma_{\mathrm{geom}} < 0.07$).

\paragraph{Not explained.}
Why Jupiter has $N = 8$ at one pole and $N = 5$ at the other (this
requires detailed vortex dynamics modeling beyond the scope of this
framework).  Why the atmospheric thermostat operates at negative
temperature (this is observed but not derived from first principles).
Why the enstrophy is minimized along the loxodromic direction
($Z''(0) > 0$ is verified computationally but not proven analytically).

\subsection{The enstrophy--energy sign structure}
\label{sec:enstrophy-discussion}

The energy--enstrophy structure at $\sigma = 0$ deserves explicit
discussion because it connects the geometric framework to Onsager's
statistical mechanics.

Along the loxodromic deformation direction in $\Cstar$, the regular
$N$-gon has
\begin{equation}\label{eq:sign-structure}
  H''(0) < 0 \quad (\text{energy maximum}), \qquad
  Z''(0) > 0 \quad (\text{enstrophy minimum}),
\end{equation}
where $H''(0) < 0$ is proven for all $N \ge 3$
(Proposition~\ref{prop:energy-max}) and $Z''(0) > 0$ is verified
computationally for $N = 3, \ldots, 9$.  The ratio $\mu = -Z''/H'' > 0$
for all $N$, placing the polygon at the \emph{negative-temperature}
state: maximum energy, minimum enstrophy.

The enstrophy minimum is \emph{local}, restricted to the one-parameter
loxodromic deformation family.  In the full $N$-dimensional
configuration space, perturbations exist that decrease $Z$ (random tests
show 62--78\% of perturbations increase $Z$, not 100\%).  The polygon is
not the global enstrophy minimizer, but it occupies the distinguished
position along the physically relevant deformation direction.

The physical picture: once a dynamical mechanism delivers the system to
$\sigma \approx 0$, the energy maximum provides no spontaneous collapse
toward a dispersed state (that would require decreasing $H$, violating
energy conservation in the large-scale dynamics).  The enstrophy minimum
provides no force away from the polygon along the loxodromic direction.
Branch disconnection (within WKB) provides no perturbative exit from
$\sigma = 0$ at leading order.  The polygon is held in place by the
convergence of conservation laws, variational structure, and eigenvalue
branch structure.

\subsection{Gaps in the five-step argument}

Three logical gaps connect the five steps and should be stated explicitly.

\paragraph{The QGPV--vortex bridge.}
Steps 1--2 operate on the QGPV (a PDE for a continuous fluid); Steps 3--5
operate on the Thomson point vortex system (discrete vortices with
logarithmic interaction).  The universality claim requires that these two
descriptions apply to the same physical system.  For Jupiter's cyclone
rings (discrete storms), the point vortex model applies directly.  For
Saturn's hexagonal jet meander (a continuous flow), the identification
of the jet with a vortex ring is a modeling assumption, not a derivation.
The bridge between the two descriptions---when and why a continuous jet
can be represented by $N$ discrete vortices---is part of the broader
vortex-dynamics literature \citep{Aref2003} but is not established within
this paper.

\paragraph{Onsager applicability.}
Step 4 invokes Onsager's (1949) negative-temperature condensation for
point vortices.  Applying this to planetary atmospheres requires three
assumptions not proven here: (a) the large-scale flow is well-described
by a point vortex system, (b) the system is in statistical equilibrium
at negative temperature, and (c) the energy maximum of the point vortex
ring coincides with the polygonal state of the continuous flow.  These
are physically reasonable for Jupiter's cyclone rings (discrete vortices
in approximate statistical equilibrium) but less directly applicable to
Saturn's jet meander.

\paragraph{Amplitude uncertainty.}
The $\varepsilon_{\mathrm{pred}} = 0.112$ matching the Cassini
observation to $0.4\%$ overstates the precision.  The matching constant
$C = 0.797$ is computed from a Gaussian jet approximation with 9.8\%
RMS fit error.  Propagating this profile uncertainty through the
formula gives $\varepsilon_{\mathrm{pred}} = 0.11 \pm 0.11$---the
true uncertainty is of order 100\%, not $0.4\%$.  The match is within
one standard deviation of the Gaussian approximation, which is
consistent but not sharp.  Furthermore, $U^*$ appears in both the
numerator ($\delta U = U_{\mathrm{obs}} - U^*$) and denominator of
$\varepsilon = C \cdot \delta U / U^*$, making the prediction
sensitive to the \emph{internal consistency} of the Gaussian
approximation rather than to the actual wind profile.

\subsection{Limitations}

\begin{itemize}
  \item \textbf{Barotropic approximation.}  The 12\% discrepancy between
    $U^*$ and $U_{\mathrm{obs}}$ for Saturn reflects the neglect of
    vertical structure.

  \item \textbf{Digitized data.}  Cassini wind profiles and Juno cyclone
    positions are digitized from published figures.  Quantitative matches
    should be verified against PDS archival data.

  \item \textbf{Limited sample.}  Three planetary configurations from two
    planets.  Uranus and Neptune observations would strengthen or falsify
    the universality claim.

  \item \textbf{Statistical, not deterministic.}  The drive toward the
    polygon is statistical (most probable macrostate), not deterministic
    (Hamiltonian trajectory).  The atmospheric thermostat is observed but
    not derived from first principles.

  \item \textbf{Local enstrophy minimum.}  The sign structure
    \eqref{eq:sign-structure} holds along the loxodromic direction, not
    globally in configuration space.
\end{itemize}

\subsection{Future directions}

\begin{enumerate}
  \item Verify all predictions against PDS (Cassini) and JIRAM (Juno)
    archival data.
  \item Compute the Onsager partition function on $\Cstar$ to confirm
    that the negative-temperature equilibrium peaks at $\sigma = 0$.
  \item Extend to baroclinic QGPV to resolve the 12\% $U^*$ discrepancy.
  \item Apply the framework to Uranus and Neptune when high-resolution
    polar data become available.
  \item Derive the $N$-selection mechanism for Jupiter: does the central
    vortex strength $\kappa_0/\kappa$ determine $N$ through the Thomson
    critical ratio hierarchy?
  \item Test on rotating-tank experiments \citep{BarbosaAguiar2010}, where
    the thermostat (viscous dissipation) is known exactly.
\end{enumerate}
